\section{Definitions and Preliminaries}
\label{sec:definitions}

\subsection{Data Model and Sample}

\begin{definition}[Sample]
A \emph{sample} $S$ of size $n$ is an ordered tuple
\[
S = \big((x_1, y_1), \ldots, (x_n, y_n)\big)
\]
where each $x_i \in X$ lies in a finite metric space $(X, d)$ and each label $y_i \in \{0, 1\}$.
\end{definition}

The order of the tuple is semantically relevant only for deterministic tie-breaking. Two identical labeled points appearing at different indices are treated as distinct sample occurrences.

\begin{itemize}
    \item \textbf{Duplicate points are allowed.}
    \item \textbf{Conflicting labels at the same point are allowed.}
\end{itemize}

Any theorem or experiment that excludes duplicates or conflicts must state this explicitly.

\begin{definition}[Finite Metric Space]
A finite metric space is a finite set $X$ together with a distance function
\[
d : X \times X \to \mathbb{R}_{\ge 0}
\]
satisfying for all $x, x', z \in X$:
\begin{enumerate}
    \item $d(x, x) = 0$;
    \item $d(x, x') > 0$ for $x \ne x'$;
    \item $d(x, x') = d(x', x)$;
    \item $d(x, z) \le d(x, x') + d(x', z)$.
\end{enumerate}
\end{definition}

\begin{definition}[Graph Shortest-Path Metric]
For a finite, simple, connected, unweighted, undirected graph $G = (V, E)$, the graph metric is
\[
d_G(u, v) = \text{length of a shortest path from } u \text{ to } v.
\]
\end{definition}

Disconnected graphs are rejected in code paths that require a metric.

\subsection{Prediction Rule}

The query point $x$ may be any element of $X$, including a point that appears in the training sample.

\begin{note}[Deterministic Tie-Breaking]
Neighbor ordering is determined lexicographically by:
\begin{enumerate}
    \item smaller distance to the query point;
    \item smaller sample index in the ordered tuple.
\end{enumerate}
If class votes tie after the first $k$ ordered neighbors are selected, the predicted label is resolved by fixed label order $0 \prec 1$. Ties are always broken in favor of label $0$.
\end{note}

\begin{definition}[$k$-NN Prediction]
Given $k \in \mathbb{N}$ with $1 \le k \le n$, the deterministic $k$-NN predictor $h_S^{(k)}(x)$ is obtained by:
\begin{enumerate}
    \item ordering the sample occurrences by the deterministic neighbor rule above;
    \item selecting the first $k$ occurrences;
    \item taking the majority vote of their labels;
    \item resolving label-vote ties by $0 \prec 1$.
\end{enumerate}
\end{definition}

\subsection{Loss Function}

\begin{definition}[Binary 0-1 Loss]
For a hypothesis $h$ and a labeled point $(x, y) \in X \times \{0,1\}$, the binary classification loss is
\[
\ell(h, (x, y)) = \mathbf{1}\{h(x) \ne y\}.
\]
\end{definition}

\subsection{Sample Perturbations}

Let $S = \big((x_1, y_1), \ldots, (x_n, y_n)\big)$.

\begin{definition}[Delete-One]
For index $i \in \{1, \ldots, n\}$, define
\[
S^{-i} = \big((x_1, y_1), \ldots, (x_{i-1}, y_{i-1}), (x_{i+1}, y_{i+1}), \ldots, (x_n, y_n)\big).
\]
\end{definition}

\begin{definition}[Add-One]
For any labeled point $z = (x, y) \in X \times \{0,1\}$, define
\[
S \oplus z
\]
as the ordered tuple obtained by appending $z$ at the end.
\end{definition}

\begin{definition}[Replace-One]
For index $i$ and labeled point $z = (x', y') \in X \times \{0,1\}$, define
\[
S^{i \leftarrow z}
\]
as the ordered tuple obtained by replacing the $i$-th occurrence of $S$ with $z$.
\end{definition}

In v0.1 worst-case stability definitions, the replacement point ranges over all labeled points in $X \times \{0,1\}$, not over a data-generating distribution.

\subsection{Stability Notions}

Unless otherwise stated, all pointwise notions are evaluated with fixed $k$, fixed metric space, fixed tie-breaking rule, and fixed loss defined above.

\begin{note}
Pointwise stability statements concern a fixed training sample, a fixed perturbation, and a fixed evaluation pair $(x, y)$. Expected stability refers to the expectation under an explicitly specified data-generating distribution. Uniform stability is the supremum of the corresponding indicator over all valid ordered training tuples, all allowed perturbation choices, and all evaluation pairs.
\end{note}

\begin{definition}[Pointwise Delete-One Stability Indicator]
For index $i$ and query-label pair $(x, y)$, define
\[
\Delta_{\mathrm{del}}(S, i, x, y)
=
\left| \ell(h_S^{(k)}, (x, y)) - \ell(h_{S^{-i}}^{(k')}, (x, y)) \right|,
\]
where $k' = \min(k, n-1)$ when deletion reduces sample size below $k$.
\end{definition}

\begin{definition}[Pointwise Replace-One Stability Indicator]
For index $i$, replacement $z$, and query-label pair $(x, y)$, define
\[
\Delta_{\mathrm{rep}}(S, i, z, x, y)
=
\left| \ell(h_S^{(k)}, (x, y)) - \ell(h_{S^{i \leftarrow z}}^{(k)}, (x, y)) \right|.
\]
\end{definition}

\begin{definition}[Pointwise Add-One Stability Indicator]
For added labeled point $z$ and query-label pair $(x, y)$, define
\[
\Delta_{\mathrm{add}}(S, z, x, y)
=
\left| \ell(h_S^{(k)}, (x, y)) - \ell(h_{S \oplus z}^{(k)}, (x, y)) \right|.
\]
\end{definition}

\begin{definition}[LOO / CVloo Stability]
For each index $i$, define
\[
\Delta_{\mathrm{loo}}(S, i)
=
\left| \ell(h_{S^{-i}}^{(k')}, (x_i, y_i)) - \ell(h_S^{(k)}, (x_i, y_i)) \right|.
\]
\end{definition}

LOO evaluates at the deleted sample occurrence itself. This definition is occurrence-based: if the same point appears twice, deleting one copy and evaluating at that occurrence is distinct from deleting the other.

For LOO, the uniform stability supremum is over ordered tuples $S$ and indices $i$, with the evaluation pair fixed to the deleted occurrence.

\subsection{Consistency Note}

The consistency notion used in paper-level discussion is classical distributional consistency for $k$-NN classification. It is not a finite-sequence surrogate notion. Any finite-metric experiment provides evidence about worst-case stability behavior; it does not by itself constitute a consistency statement.
